\subsection{The coefficient-$p$ half at derived length three}

Throughout this subsection, $A$ is the complete normal two-dimensional
local domain of equal characteristic $p>0$ with algebraically closed
residue field used in F8--F11, $U$ is its punctured spectrum, and
all resolutions are fixed SNC resolutions. No mixed-characteristic
assertion is made here. The proof passed an independent root audit.

\begin{lemma}[F12a: normal $p$-subgroups on simple characteristic-$p$ modules]
Let $\Gamma$ be a finite group and let $V$ be a nonzero simple
$\mathbf F_p[\Gamma]$-module. Every normal $p$-subgroup of
$\Gamma$ acts trivially on $V$. Consequently, if $\Gamma$ is
metabelian and its action on $V$ is faithful, then $\Gamma'$ has
order prime to $p$.
\end{lemma}
\begin{proof}
Every nonzero finite-dimensional representation of a finite
$p$-group $P$ over $\mathbf F_p$ has nonzero invariants. Here is an
inductive proof. A central subgroup $C_p\subset P$ exists; its
generator $u$ satisfies $(u-1)^p=0$, so the kernel of $u-1$ is
nonzero and $P$-stable. The quotient group $P/C_p$ acts on it,
and induction gives a nonzero $P$-fixed vector. If $P$ is normal
in $\Gamma$, the invariant space $V^P$ is $\Gamma$-stable, so
simplicity gives $V^P=V$.
For metabelian $\Gamma$, the $p$-primary part of its abelian
commutator group is characteristic in $\Gamma'$ and therefore
normal in $\Gamma$. It acts trivially; faithfulness makes it
trivial.
\end{proof}

\begin{proposition}[F12: characteristic-$p$ coefficients with metabelian monodromy; Proved]
There is a constant $C_{A,2,p}$ such that every simple finite
$\mathbf F_p$-local system $M$ on $U$ whose monodromy group is
metabelian satisfies
\[
                    \dim_{\mathbf F_p}H^1(U,M)
                    \leq C_{A,2,p}\operatorname{rank}M.
 \tag{F12.1}
\]
The system itself may be wildly ramified on the original resolution.
\end{proposition}
\begin{proof}
Use once and for all the finite connected abelian cover
$U_0\to U$ and the fixed resolution $X_0$ supplied by F10b.
Write $F=\operatorname{Gal}(U_0/U)$ and let
$\rho:\pi_1(U)\twoheadrightarrow\Gamma\subset\operatorname{GL}(M)$ be the
faithful monodromy image. It is metabelian by hypothesis.
Lemma F12a makes $\Gamma'$ a prime-to-$p$ group.

The finite quotient $\pi_1(U)\to\Gamma^{\mathrm{ab}}$ is an
abelian local system. Its pullback to $U_0$ extends to $X_0$,
by the uniform assertion of F10b. Thus at every exceptional
DVR of $X_0$, the image under $\rho$ of boundary inertia lies
in $\Gamma'$. Wild inertia is pro-$p$, so its image in the
prime-to-$p$ group $\Gamma'$ is trivial. Therefore the entire
pulled-back system $M|_{U_0}$ is tame along the fixed exceptional
divisor of $X_0$.

No simplicity assumption on this restricted system is needed:
F8 applies to arbitrary finite $\mathbf F_p$-local systems with
tame boundary monodromy. Hence
\[
 \dim H^1(U_0,M)\leq C_{X_0}\operatorname{rank}M.
\]
Inflation--restriction bounds the kernel of restriction by
$\dim H^1(F,M^{\pi_1(U_0)})\leq d(F)\operatorname{rank}M$.
Thus $C_{A,2,p}=C_{X_0}+d(F)$ works. In particular this argument
does not assume that $[U_0:U]$ is prime to $p$ or that restriction
preserves simplicity.
\end{proof}

\begin{corollary}[F12b: a pro-$p$ second derived subgroup; Proved]
Every quotient $G$ of $\pi_1(U)$ whose closed second derived
subgroup $G''$ is pro-$p$ is topologically finitely generated,
with a bound depending only on $A$.
\end{corollary}
\begin{proof}
The quotient $\Gamma=G/G''$ is metabelian, so F11 bounds its
number of generators uniformly. Every finite simple
$\mathbf F_p[G]$-module kills the normal pro-$p$ group $G''$ by
F12a applied to its finite image. Its faithful monodromy group
is therefore metabelian, so inflation and F12 give the uniform
$H^1$ bound. The relative criterion F7a proves the assertion.
\end{proof}

\paragraph{Precise remaining derived-length-three gate.}
For a finite group of derived length at most three, a minimal
normal subgroup inside its second derived group is an elementary
abelian $\ell$-group. Its conjugation representation is simple
and factors through a metabelian group, because the second derived
group is abelian. F12 supplies the needed uniform estimate when
$\ell=p$. The cases $\ell\ne p$, with possibly wild metabelian
coefficient monodromy, remain open here. F12b is a consequence
for a specified class, not a claim that every second derived group
is pro-$p$ or that bounded-exponent inertia inside a finitely
generated metabelian group is finite.
